\section{Техническая реализация программных модулей}

В этой главе мы разберем исходный код инструментов, которые используются для создания карт и навигации. Всё здесь --- обычный Python на компьютере: дрон и симулятор понадобятся только в последнем задании, чтобы получить журнал настоящего полёта.

\subsection{Генератор карт (Marker Map Builder)}
Скрипт \texttt{19\_map\_creation\_example3\_builder.py} (глава «Создание навигационных карт») позволяет создавать цифровые карты полигонов и экспортировать их в форматы JSON, XML, KML, YAML, PDF и PNG.

\subsubsection{Структура класса Marker}
Каждый маркер описывается классом \texttt{Marker}, который хранит:
\begin{itemize}
    \item \texttt{id}: Уникальный номер маркера.
    \item \texttt{size}: Физический размер стороны в метрах.
    \item \texttt{x, y, z}: Локальные координаты относительно начала координат карты.
    \item \texttt{yaw}: Угол поворота (ориентация), градусы.
    \item \texttt{dictionary}: Тип словаря (например, DICT\_6X6\_250).
\end{itemize}

\begin{codefiletitle}[python]{Класс Marker}{python/advanced_navigation/marker_class_code.py}
\end{codefiletitle}

\begin{taskbox}[Свой парсер]
\emph{Компьютер.} Напишите функцию \texttt{load\_json(filename)}, которая считывает файл карты (формат метода \verb|save_json|) и возвращает список объектов класса \texttt{Marker}. Сохраните сетку $10 \times 10$ из конструктора, загрузите её обратно и сравните.
\expected{после загрузки 100 маркеров; у каждого \texttt{id}, \texttt{size}, \texttt{x}, \texttt{y}, \texttt{z}, \texttt{yaw} и \texttt{dictionary} совпадают с исходными (проверка \texttt{assert} в цикле проходит без ошибок).}
\end{taskbox}

\subsubsection{Экспорт в XML}
Метод \texttt{save\_xml} генерирует структуру, понятную большинству промышленных систем. Мы используем f-строки для формирования XML-тегов.

\begin{codefiletitle}[python]{Экспорт в XML}{python/advanced_navigation/save_xml_code.py}
\end{codefiletitle}

У f-строк есть слабое место: они не экранируют специальные символы. Если в названии карты встретится \verb|&| или \verb|<|, получится файл, который ни одна программа не прочитает. Надёжнее строить XML модулем \verb|xml.etree.ElementTree| --- он сам экранирует текст.

\begin{taskbox}[XML без сюрпризов]
\emph{Компьютер.} Назовите карту \verb|"Склад & цех"| и сохраните её методом \verb|save_xml|. Попробуйте прочитать файл: \verb|xml.etree.ElementTree.parse(filename)|. Затем напишите \verb|save_xml_etree(self, filename)|, которая строит то же дерево через \verb|ET.Element| и \verb|ET.SubElement| и сохраняет \verb|ET.ElementTree(root).write(filename, encoding="utf-8", xml_declaration=True)|.
\expected{файл из f-строк не читается (\texttt{ParseError}); файл из \texttt{ElementTree} читается, а \texttt{root.find("name").text} возвращает исходное название \texttt{Склад \& цех}.}
\end{taskbox}

\subsubsection{Экспорт в KML для Google Earth}
Для отображения маркеров на глобусе Google Earth используется формат KML. Главная сложность здесь — перевод локальных метров $(x, y)$ в глобальные градусы широты и долготы. Её решает метод \verb|local_to_gps| конструктора.

Мы используем приближенную формулу (для малых расстояний):
\begin{itemize}
    \item 1 градус широты $\approx 111132$ метра.
    \item 1 градус долготы $\approx 111132 \cdot \cos(\text{широта})$ метров.
\end{itemize}
Кроме того, нужно знать, куда смотрел дрон на старте. Конструктор считает, что на север: тогда X Пионера --- на север (широта растёт), а Y, направленный влево, --- на запад (долгота \emph{уменьшается}). В самом файле KML координаты записываются в порядке «долгота, широта, высота».

\begin{warnbox}[Точность GPS]
Данная формула имеет погрешность. Для профессиональной геодезии следует использовать библиотеку \texttt{pyproj} и проекцию UTM.
\end{warnbox}

\begin{taskbox}[Проверка пересчёта]
\emph{Компьютер.} Для карты с началом в точке $(59.9343;\ 30.3351)$ вызовите \verb|local_to_gps(100, 0)| и \verb|local_to_gps(0, 100)|. Затем откройте KML-файл в Google Earth или QGIS и линейкой программы измерьте расстояние между маркерами 0 и 9 сетки $10 \times 10$.
\expected{$(59.93520;\ 30.33510)$ и $(59.93430;\ 30.33330)$ с точностью до пятого знака; расстояние между маркерами 0 и 9 --- 4,5~м с ошибкой не больше нескольких сантиметров.}
\end{taskbox}

\subsection{Модуль визуализации}
Для проверки правильности карты полезно отрисовать её на 2D-плоскости перед полетом. Библиотека \texttt{matplotlib} позволяет построить схему расположения маркеров. Функция ниже написана как метод класса \verb|MarkerMap| (отсюда параметр \verb|self|).

\begin{codefiletitle}[python]{Визуализация карты}{python/advanced_navigation/plot_map_code.py}
\end{codefiletitle}

\begin{taskbox}[Трек полёта поверх карты]
\emph{Симулятор, затем компьютер.} В симуляторе выполните секторный поиск метки 12 (миссия Python №\,101) и после каждой команды печатайте строку \verb|track;x;y;z| из \verb|get_local_position_lps()|, а при каждом чтении кадра --- \verb|seen;id1,id2| со списком прочитанных ID. Скопируйте журнал из терминала в файл \verb|track.txt|. На компьютере постройте карту с маркером 12 и нанесите на неё трек: строки \verb|track| --- линией, места, где была прочитана метка 12, --- крестиками.
\expected{рисунок, на котором линия трека проходит через три сектора с шагом 4~м по $Y$, а крестики стоят в том секторе, где метка лежит на карте; после находки трек идёт обратно к старту.}
\mission{Python}{101}
\end{taskbox}
